\documentclass[11pt,a4paper]{article}
\usepackage[T1]{fontenc}
\usepackage[utf8]{inputenc}
\usepackage{lmodern}
\usepackage{amsmath,amssymb,amsthm,mathtools,mathrsfs}
\usepackage{graphicx,float,xcolor,multicol}
\usepackage[shortlabels]{enumitem}
\usepackage[margin=27mm]{geometry}
\usepackage{microtype,needspace,placeins}
\usepackage{tikz,tikz-cd}
\usetikzlibrary{arrows.meta,calc,positioning,patterns,decorations.pathreplacing}
\usepackage[colorlinks=true,linkcolor=teal,urlcolor=teal,citecolor=teal]{hyperref}
\setlength{\parindent}{0pt}
\setlength{\parskip}{.6em}
\setcounter{secnumdepth}{0}
\allowdisplaybreaks
\raggedbottom
\providecommand{\cross}{\times}
\DeclareUnicodeCharacter{2020}{\ensuremath{\dagger}}
\DeclareMathOperator{\supp}{supp}
\DeclareMathOperator*{\esssup}{ess\,sup}
\providecommand{\chapter}{\section}
\newenvironment{tool}[2]{\subsection*{Tool #1 --- #2}}{}
\newcommand{\ToolRef}[1]{Tool #1}
\newcommand{\FigureTag}[2]{\textit{Figure #1}}
\newcommand{\Result}[1]{\subsection*{#1}}
\newcommand{\Application}[1]{\subsubsection*{Application\if\relax\detokenize{#1}\relax\else\space--- #1\fi}}
\newcommand{\ProofHeading}{\par\textit{Proof.}\space}
\newcommand{\ProofOf}[1]{\par\textit{Proof of #1.}\space}
\newcommand{\RemarkHeading}{\par\textbf{Remark.}\space}
\newcommand{\NoteHeading}{\par\textbf{Note.}\space}
\newcommand{\ExerciseHeading}{\par\textbf{Exercise.}\space}

% Rendering support only; mathematical text is in each independent post.
\usepackage{booktabs,array,longtable}
\definecolor{HandoutBlue}{HTML}{2F6690}
\definecolor{HandoutNavy}{HTML}{17365D}
\newtheorem{theorem}{Theorem}
\newtheorem{lemma}[theorem]{Lemma}
\newtheorem{proposition}[theorem]{Proposition}
\newtheorem{corollary}[theorem]{Corollary}
\newtheorem{conjecture}[theorem]{Conjecture}
\newtheorem{definition}{Definition}
\newtheorem{example}{Example}
\newtheorem{namedtoolcounter}{Tool}
\newenvironment{namedtool}[1]{\begin{namedtoolcounter}[#1]}{\end{namedtoolcounter}}
\newtheorem*{claim}{Claim}
\newtheorem*{remark}{Remark}
\newtheorem*{exercise}{Exercise}
\newtheorem*{question}{Question}
\newtheorem*{observation}{Observation}
\newcommand{\SourceChapter}[2]{\section*{#2}}
\newcommand{\Lecture}[2]{\section*{Lecture #1: #2}}
\newcommand{\unclear}[1]{\textit{[unclear: #1]}}
\newcommand{\sourcegap}{\textit{[unfinished in the source]}}
\DeclareMathOperator{\dist}{dist}
\DeclareMathOperator{\id}{id}
\DeclareMathOperator{\Hol}{Hol}
\DeclareMathOperator{\Per}{Per}
\DeclareMathOperator{\Fix}{Fix}
\DeclareMathOperator{\diam}{diam}
\DeclareMathOperator{\Var}{Var}
\DeclareMathOperator{\Leb}{Leb}
\DeclareMathOperator{\Hom}{Hom}
\DeclareMathOperator{\End}{End}
\DeclareMathOperator{\Ann}{Ann}
\DeclareMathOperator{\Spec}{Spec}
\DeclareMathOperator{\Max}{Max}
\DeclareMathOperator{\Ob}{Ob}
\DeclareMathOperator{\Tor}{Tor}
\DeclareMathOperator{\Ass}{Ass}
\DeclareMathOperator{\Mor}{Mor}
\DeclareMathOperator{\Fun}{Fun}
\DeclareMathOperator{\Nat}{Nat}
\DeclareMathOperator{\Cone}{Cone}
\DeclareMathOperator{\MCG}{MCG}
\DeclareMathOperator{\Aut}{Aut}
\DeclareMathOperator{\Deck}{Deck}
\DeclareMathOperator{\SL}{SL}
\DeclareMathOperator{\PSL}{PSL}
\DeclareMathOperator{\Area}{Area}
\DeclareMathOperator{\im}{im}
\DeclareMathOperator{\PSh}{PSh}
\newcommand{\CRing}{\mathsf{CRings}}
\newcommand{\AMod}{A\text{-}\mathsf{Mod}}
\newcommand{\Ab}{\mathsf{Ab}}
\newcommand{\Z}{\mathbb Z}
\newcommand{\Q}{\mathbb Q}
\newcommand{\R}{\mathbb R}
\newcommand{\C}{\mathbb C}
\newcommand{\N}{\mathbb N}
\newcommand{\kk}{\mathbb k}
\renewcommand{\aa}{\mathfrak a}
\newcommand{\bb}{\mathfrak b}
\newcommand{\pp}{\mathfrak p}
\newcommand{\qq}{\mathfrak q}
\newcommand{\mm}{\mathfrak m}
\newcommand{\nn}{\mathfrak n}
\newcommand{\rad}{\mathrm r}
\newcommand{\cat}[1]{\mathsf{#1}}
\setcounter{secnumdepth}{0}
\allowdisplaybreaks

\graphicspath{{figures/lemma-book/}}
\begin{document}
\section*{An Inequality Toolkit: Jensen, Muirhead, and Schur}
% BLOG-CONTENT-BEGIN
Prove that $\dfrac a{1+bc}+\dfrac b{1+ca}+\dfrac c{1+ab}\leq2$ if
$0\leq a,b,c\leq1$.
For fixed $b,c$, the expression
\[
\frac a{1+bc}+\frac b{1+ca}+\frac c{1+ab}
\]
is a convex function of $a$ on $[0,1]$. Hence its maximum occurs at
$a=0$ or $a=1$. Repeating this argument in $b$ and $c$, its maximum on
$[0,1]^3$ occurs at a vertex. Checking the eight vertices gives the
upper bound $2$.

Here is the original route in the notebook. Without loss of generality, let $a\geq b\geq c$.
Since $(1-a)(1-b)\geq0$, we have $1+ab\geq a+b$, so $1+2ab\geq a+b$.
Now $a+b+c\leq1+1+2ab=2(1+ab)$, and hence the notebook continues with
\[
\frac a{1+bc}+\frac b{1+ca}+\frac c{1+ab}
\leq\frac a{1+ab}+\frac b{1+ab}+\frac c{1+ab}\leq2.
\]
The first displayed inequality has the denominator comparison in the wrong direction; for example, $a=1$, $b=\tfrac12$, $c=0$ makes its left side $3/2$ and its middle expression $1$. The convexity argument above corrects this step without removing the original attempt.

\setcounter{theorem}{14}\begin{lemma}[Cauchy's inequality application]
If $x,y,z>1$, then
$\dfrac{x^4}{(y-1)^2}+\dfrac{y^4}{(z-1)^2}+\dfrac{z^4}{(x-1)^2}\geq48$.
\end{lemma}
Cauchy's inequality reduces the desired assertion to
\[
\left(\frac{x^2}{y-1}+\frac{y^2}{z-1}+\frac{z^2}{x-1}\right)^2\geq144,
\quad\text{and then to}\quad
\frac{x^2}{y-1}+\frac{y^2}{z-1}+\frac{z^2}{x-1}\geq12.
\]
It remains to show
$\dfrac{(x+y+z)^2}{x+y+z-3}\geq12$.
Writing $a=x+y+z$, this becomes
$a^2\geq12a-36$, or $(a-6)^2\geq0$.

% Source page 7 - left
\setcounter{theorem}{21}\begin{lemma}[Catalan's identity]
$1-\frac12+\frac13-\frac14+\cdots+\frac1{2n-1}-\frac1{2n}
=\sum_{k=1}^n\frac1{n+k}$.
\end{lemma}
\paragraph{Application.}
For every integer $n\geq2$, prove that
\[
\frac1{n+1}\left(1+\frac13+\frac15+\cdots+\frac1{2n-1}\right)
>\frac1n\left(\frac12+\frac14+\cdots+\frac1{2n}\right).
\]
Multiplying through, the desired inequality becomes
\[
1+\frac13+\cdots+\frac1{2n-1}
>\left(1+\frac1n\right)\left(\frac12+\frac14+\cdots+\frac1{2n}\right),
\]
or
\[
1-\frac12+\frac13-\frac14+\cdots+\frac1{2n-1}-\frac1{2n}
>\frac1{2n}+\frac1{4n}+\frac1{6n}+\cdots+\frac1{2n^2}.
\]
By the identity, it is enough to compare
$\frac1{n+1}+\cdots+\frac1{2n}$ with $\frac1{2n}+\frac1{4n}+\cdots+\frac1{2n^2}$.
Add the inequalities
$\frac1{2n}\geq\frac1{2n}$, $\frac1{2n-1}>\frac1{4n}$, $\ldots$,
$\frac1{n+1}>\frac1{2n^2}$ to finish.

\setcounter{theorem}{22}\begin{lemma}[Homogenization]
Homogenization rewrites an inequality so that all terms have the same degree.
\end{lemma}
\paragraph{Application.}
Let $a,b,c>0$ and $abc=1$. Prove that $a+b+c\leq a^2+b^2+c^2$.
% Source page 7 - right
There are expressions of three different degrees, $3$, $2$, and $1$.
To homogenize the inequality, multiply by $\sqrt[3]{abc}$:
\[
\sqrt[3]{abc}(a+b+c)\leq a^2+b^2+c^2,
\quad\text{or}\quad
a^{4/3}b^{1/3}c^{1/3}+a^{1/3}b^{4/3}c^{1/3}+a^{1/3}b^{1/3}c^{4/3}
\leq a^2+b^2+c^2.
\]
Since $(4/3,1/3,1/3)\prec(2,0,0)$, Muirhead's theorem gives
$\sum_{\mathrm{cyc}}a^{4/3}b^{1/3}c^{1/3}\leq\sum_{\mathrm{cyc}}a^2$.

\setcounter{theorem}{23}\begin{lemma}[Second-derivative test]
The best test for convexity or concavity is the second-derivative test:
if $f''(x)\geq0$, then $f$ is convex, while $f''(x)\leq0$ implies that $f$ is concave.
\end{lemma}
This technique, combined with Jensen's inequality or a maxima--minima
argument, yields many beautiful and elegant solutions.

\setcounter{theorem}{24}\begin{lemma}[Jensen's inequality]
Let $f:I\to\mathbb R$ be convex. For any $x_1,\ldots,x_n\in I$ and nonnegative reals $\omega_1,\ldots,\omega_n$, not all zero,
\[
\omega_1f(x_1)+\cdots+\omega_nf(x_n)
\geq\left(\sum_{i=1}^n\omega_i\right)
 f\left(\frac{\sum_i\omega_ix_i}{\sum_i\omega_i}\right).
\]
If $f$ is concave, reverse the sign.
\end{lemma}
\paragraph{Application: IMO 1995, problem 2.}
If $a,b,c>0$ and $abc=1$, prove that
\[
\frac1{a^3(b+c)}+\frac1{b^3(c+a)}+\frac1{c^3(a+b)}\geq\frac32.
\]
% Source page 8 - left
Substitute $x=1/a$, $y=1/b$, $z=1/c$. For $f(x)=1/x$, we have $f''(x)=2/x^3>0$, so $f$ is convex.
Now Jensen's inequality gives
\begin{align*}
\frac1{a^3(b+c)}+\frac1{b^3(c+a)}+\frac1{c^3(a+b)}
&=\frac{x^2}{y+z}+\frac{y^2}{z+x}+\frac{z^2}{x+y}\\
&=xf\left(\frac{y+z}{x}\right)+yf\left(\frac{z+x}{y}\right)+zf\left(\frac{x+y}{z}\right)\\
&\geq(x+y+z)f\left(\frac{(y+z)+(z+x)+(x+y)}{x+y+z}\right)\\
&=\frac{x+y+z}{2}\geq\frac{3\sqrt[3]{xyz}}2
=\frac{3\sqrt[3]{1/(abc)}}2=\frac32.
\end{align*}

\setcounter{theorem}{25}\begin{lemma}[Muirhead and majorization]
Muirhead's theorem and majorization techniques are analogous.
If $(a_1,\ldots,a_n)\succ(b_1,\ldots,b_n)$, then for positive $x_1,\ldots,x_n$,
\[
\sum_{\mathrm{sym}}x_1^{a_1}\cdots x_n^{a_n}
\geq\sum_{\mathrm{sym}}x_1^{b_1}\cdots x_n^{b_n}.
\]
For majorization, let $f:I\to\mathbb R$ be convex and suppose $(x_n)\succ(y_n)$.
Then $\sum_{i=1}^n f(x_i)\geq\sum_{i=1}^n f(y_i)$ (Karamata's inequality).
\end{lemma}
\emph{For a concave function, reverse the inequality.}
Here are two applications of these ideas.
% Source page 8 - right
Prove that, for positive reals $a,b,c$,
$(a^2b+b^2c+c^2a)(ab^2+bc^2+ca^2)\geq9a^2b^2c^2$.
Expanding the product gives
\begin{align*}
(a^2b+b^2c+c^2a)(ab^2+bc^2+ca^2)
={}&a^3b^3+b^3c^3+c^3a^3+3a^2b^2c^2\\
&+a^4bc+ab^4c+abc^4.
\end{align*}
Notice $(3,3,0)\succ(2,2,2)$ and $(4,1,1)\succ(2,2,2)$.
Thus Muirhead gives
\[
\sum_{\mathrm{sym}}a^3b^3c^0\geq\sum_{\mathrm{sym}}a^2b^2c^2,
\qquad \sum_{\mathrm{sym}}a^4bc\geq\sum_{\mathrm{sym}}a^2b^2c^2.
\]
The conclusion follows.

\setcounter{theorem}{26}\begin{lemma}[Chebyshev's theorem]
Let $a_1\leq a_2\leq\cdots\leq a_n$ and $b_1\leq b_2\leq\cdots\leq b_n$ be nondecreasing real sequences. Then
\[
\frac1n\sum_{i=1}^n a_ib_i
\geq\left(\frac1n\sum_{i=1}^n a_i\right)
\left(\frac1n\sum_{i=1}^n b_i\right)
\geq\frac1n\sum_{i=1}^n a_ib_{n-(i-1)}.
\]
\end{lemma}
When this is used, its validity in the particular problem must be closely inspected.
\paragraph{Application.}
Let us resolve the same IMO 1995, problem 2, mentioned above.
Apply the same substitution as before, with $x\geq y\geq z$. Then
$\dfrac{x}{y+z}\geq\dfrac{y}{z+x}\geq\dfrac{z}{x+y}$.
Chebyshev's inequality gives
\[
x\frac{x}{y+z}+y\frac{y}{z+x}+z\frac{z}{x+y}
\geq\frac{x+y+z}{3}\left(\frac{x}{y+z}+\frac{y}{z+x}+\frac{z}{x+y}\right).
\]
We know $x+y+z\geq3$ and, by Nesbitt's inequality,
$\dfrac{x}{y+z}+\dfrac{y}{z+x}+\dfrac{z}{x+y}\geq\dfrac32$.
The conclusion follows.

\setcounter{theorem}{52}\begin{lemma}[Sum-of-squares strategy]\label{lemma:lb1-58}
If none of the standard techniques proves an inequality, try expressing it as a sum of squares greater than or equal to zero.
\end{lemma}
\paragraph{Application: V. C\^{i}rtoaje.}
For real $a,b,c$, prove $(a^2+b^2+c^2)^2\geq3(a^3b+b^3c+c^3a)$.
The expression has the following sum-of-squares decomposition.
% Source page 24 - right
\[
(a^2+b^2+c^2)^2-3(a^3b+b^3c+c^3a)
=\frac12\sum_{\mathrm{cyc}}(a^2-2ab+bc-c^2+ca)^2\geq0.
\]
This completes the proof.

\setcounter{theorem}{59}\begin{lemma}[Schur's inequality]
Define the normalized symmetric sum
\[
T(a_1,\ldots,a_n)=\frac1{n!}\sum_{\sigma\in S_n}
x_1^{a_{\sigma(1)}}\cdots x_n^{a_{\sigma(n)}}.
\]
Here the variables $x_1,\ldots,x_n$ are positive. For
$\alpha\in\mathbb R$ and $\beta>0$, the following inequalities hold:
\begin{align*}
T(\alpha+2\beta,0,0)+T(\alpha,\beta,\beta)
&\geq T(\alpha+\beta,\beta,0)+T(\alpha+\beta,0,\beta),&&\text{first form},\\
T(\alpha+\beta,\alpha+\beta,0)+T(2\alpha,\beta,\beta)
&\geq2T(\alpha+\beta,\beta,\alpha),&&\text{second form}.
\end{align*}
\end{lemma}
\paragraph{Applications.}
Prove that
\[
a^4+b^4+c^4+a^2bc+b^2ca+c^2ab\geq2(a^3b+b^3c+c^3a).
\]
First realize that standard Muirhead fails; with Schur, the proposed
solution is two lines:
\[
T(2+2\cdot1,0,0)+T(2,1,1)\geq2T(3,0,0),
\]
which is obtained by setting $\alpha=2$, $\beta=1$ in the notes.
As written, however, the cyclic inequality is false: for $(a,b,c)=(1,4,3)$,
its two sides are $434$ and $446$. The displayed $T$-inequality also has
different total degrees on its two sides and does not encode that claim.

The corresponding valid symmetric form of Schur is the following.
Schur's first form with $\alpha=2$ and $\beta=1$ gives
\[
T(4,0,0)+T(2,1,1)\geq T(3,1,0)+T(3,0,1),
\]
or, equivalently,
\[
a^4+b^4+c^4+abc(a+b+c)\geq\sum_{\mathrm{sym}}a^3b.
\]

Prove $3(x^2y+y^2z+z^2x)(xy^2+yz^2+zx^2)\geq xyz(x+y+z)^3$.
First try all the nonsense tricks and conclude that the problem is
impenetrable. Then gather courage to expand that big bull, obtaining
\[
2T(4,1,1)+3T(3,3,0)+T(2,2,2)\geq6T(3,2,1).
\]
We have
\begin{align*}
[3,3,0]+[2,2,2]&\geq2[3,2,1],&&\text{I},\\
2[4,1,1]&\geq2[3,2,1],&&\text{II},\\
2[3,3,0]&\geq2[3,2,1],&&\text{III}.
\end{align*}
Here $[p,q,r]$ denotes the corresponding symmetric sum.
I follows from Schur's second form with $\alpha=1$, $\beta=2$.
Add I, II, and III to finish.

% Source page 27 - left
There are true villains on the IMO jury who pose problems like this one.
Prove
\[
x^4+y^4+z^4+3(x^2y^2+y^2z^2+z^2x^2)
\geq2\bigl(x^3(y+z)+y^3(z+x)+z^3(x+y)\bigr).
\]
It is homogeneous and symmetric, but none of the standard techniques works.
Apply the sum-of-squares strategy from Lemma~\ref{lemma:lb1-58}.
Write $[4,0,0]+3[2,2,0]\geq4[3,1,0]$.
Bring everything to the left and rearrange to obtain
$(x-y)^4+(y-z)^4+(z-x)^4\geq0$.

\setcounter{theorem}{60}\begin{lemma}[Square-root heuristic: Stanley Rabinowitz]
If $\sum_{\mathrm{cyc}}f(x)$ has three terms, then
$\sum_{\mathrm{cyc}}\sqrt{f(x)}\leq\sqrt{3\sum_{\mathrm{cyc}}f(x)}$.
\end{lemma}
Use the power-mean inequality: if $m\leq n$, then
\[
\left(\frac{a^m+b^m+c^m}{3}\right)^{1/m}
\leq\left(\frac{a^n+b^n+c^n}{3}\right)^{1/n}.
\]
Set $m=1/2$, $n=1$, obtaining
$\bigl(\sum\sqrt{f(x)}\bigr)^2/9\leq\sum f(x)/3$.
An alternative proof writes $3=1+1+1$ and applies Cauchy.
\setcounter{theorem}{61}\begin{lemma}[Extension of the root heuristic]
If $\sum f(x)$ has $n$ terms, then
\[
\sum_{\mathrm{cyc}}\sqrt[m]{f(x)}\leq\sqrt[m]{n^{m-1}\sum_{\mathrm{cyc}}f(x)}.
\]
This is analogous to the power-mean inequality.
\end{lemma}
% Source page 27 - right
\paragraph{Application.}
Prove $\sqrt x+\sqrt y+\sqrt z\leq\sqrt{3(x+y+z)}$.
Apply the square-root heuristic with $f(a)=a$.

\paragraph{A one-line identity, due to P. K. Hung.}
For positive $a,b,c$ with $a+b+c=1$, prove
$1/a+1/b+1/c\geq25/(1+48abc)$.
The following identity proves the claim:
\[
\frac1a+\frac1b+\frac1c-\frac{25}{1+48abc}
=\sum_{\mathrm{cyc}}\frac{(a+b-3c)^2(a-b)^2}{ab(1+48abc)}\geq0.
\]
% BLOG-CONTENT-END
\end{document}
